\begin{abstract}
A natural way to merge several low-rank task updates of one linear layer is to choose a single
rank-$r$ update that minimises the largest per-task functional error measured on calibration
activations. We study how sensitive this empirical minimax problem is to the calibration sample,
in a controlled synthetic layer-level setting with \CfgK{} tasks, input dimension \CfgDin{},
output dimension \CfgDout{}, rank budget $r=\CfgR$ and \CfgNcal{} calibration inputs per task.
We first check the optimiser. Across \STwoNFits{} minimax fits, the relative gap between a
Lagrangian lower bound, computed with an exactly solvable weighted reduced-rank regression, and
the achieved worst-task value is at most $\STwoMaxGapRel$, and an independent raw-data reference
agrees to within $\STwoMaxGDisc$. These floating-point checks argue against large optimisation
error on these instances; they are not certificates. The optimised calibration objective is
nevertheless optimistic: on fresh calibration draws the minimax solution's fitted worst-task
error (\STwoOptFitEEMinimax{} on average) understates its population value
(\STwoOptPopEEMinimax), and its population worst-task error is not lower than that of the
uniform-weight solution on average (mean paired difference \STwoDeltaEE{} over
\CfgNTeachers{} problem instances $\times$ \CfgNResamples{} nested calibration draws; the
pre-registered verdict is ambiguous). Replacing the empirical second moments by population
moments, an oracle that exists only in simulation, removes the average disadvantage, whereas
replacing only the error normalisers does not. Under our pre-registered attribution rule the
cause remains undetermined. Elementary, known perturbation bounds that make the dependence on moment error explicit hold
on every recorded fit, but at this sample size they are loose by a factor of at least
\SThreeMinBoundOverExcess. No real adapters were evaluated.
\end{abstract}
